\subsection{The finite soft function in dilogarithms}
\label{sec:comp-supp-soft-analytic}
The soft-photon calculation leaves one finite parameter integral at
fixed elastic scattering angle. It can be integrated before performing
the elastic angular integral. In this section momenta and energies are
in units of the electron mass, so $s>1$. Define
\begin{align}
 E&=\frac{s+1}{2\sqrt{s}},\qquad q=\frac{s-1}{2\sqrt{s}},
 \qquad \beta=\frac qE,\qquad t=-2q^2(1-c),\\
 z_e&=1-\frac t2,\qquad M_x=1-tx(1-x),
 \qquad v_x=\sqrt{1-M_x/E^2}.
\end{align}
Here $c$ is the cosine of the elastic photon scattering angle in the
center-of-mass frame. The function to be integrated is
\begin{equation}
 K=z_e I-\frac{\operatorname{artanh}\beta}{\beta},
 \qquad
 I=\int_0^1\frac{\operatorname{artanh}v_x}{v_xM_x}\,\dd x.
 \label{eq:comp-soft-I-definition}
\end{equation}
The ratio $\operatorname{artanh}v_x/v_x$ equals one at $v_x=0$ by continuity.
Every quantity in this integral is real in the physical angular interval.

\paragraph{Rationalizing the square root.}
For an interior angle introduce
\begin{equation}
 a=\frac{\sqrt{-t}}2,\qquad
 \kappa=\frac{\sqrt{q^2-a^2}}E,\qquad
 0<a<q,\qquad 0<\kappa<\beta<1.
\end{equation}
Symmetry about $x=1/2$ allows the substitution $x=(1+y)/2$ with
$0\leq y\leq1$. Then
\begin{equation}
 M_x=1+a^2-a^2y^2,\qquad
 v_x^2=\kappa^2+\frac{a^2y^2}{E^2},\qquad
 I=\frac1{aE}\int_\kappa^\beta
 \frac{\operatorname{artanh}v}{(1-v^2)\sqrt{v^2-\kappa^2}}\,\dd v.
\end{equation}
Set $v=\kappa\cosh u=\kappa(r+r^{-1})/2$ with $r=e^u$.
The result contains only rational functions and logarithms,
\begin{align}
 I&=\frac1{aE}\int_1^{R}
 \frac{2r}{P_+(r)P_-(r)}
 \ln\frac{P_+(r)}{P_-(r)}\,\dd r,
 \label{eq:comp-soft-rational-integral}\\
 P_+(r)&=\kappa r^2+2r+\kappa,\qquad
 P_-(r)=-\kappa r^2+2r-\kappa,\\
 R&=\frac{q+a}{\sqrt{q^2-a^2}},\qquad Y=R-1.
\end{align}
This transformation selects the positive square root because $y\geq0$
and $u\geq0$. Thus it does not require analytic continuation from another
kinematic region.

\paragraph{A real dilogarithmic primitive.}
Let $d=\sqrt{1-\kappa^2}$ and list the four roots with their signs,
\begin{align}
 (z_1,z_2,z_3,z_4)&=
 \left(-\frac{1+d}{\kappa},-\frac{\kappa}{1+d},
 \frac{\kappa}{1+d},\frac{1+d}{\kappa}\right),\\
 (\epsilon_1,\epsilon_2,\epsilon_3,\epsilon_4)&=(1,1,-1,-1),
 \qquad h_j=z_j-1,\\
 A_j&=\frac1{4-2\kappa^2(z_j^2+1)},\qquad
 L_0=\ln\frac{1+\kappa}{1-\kappa}.
\end{align}
Partial fractions and factorization give
\begin{align}
 \frac{2r}{P_+(r)P_-(r)}&=\sum_{j=1}^4\frac{A_j}{r-z_j},\\
 \ln\frac{P_+(1+y)}{P_-(1+y)}
 &=L_0+\sum_{j=1}^4\epsilon_j\ln(1-y/h_j).
\end{align}
The two negative roots and the smaller positive root have $h_j<0$.
The remaining root satisfies $z_4>R$, because $v(R)=\beta<1$
whereas $v(z_4)=1$. Consequently every factor $1-y/h_j$ is positive
for $0\leq y\leq Y$. This proves that all logarithms in the transformed
integral are real and that no pole lies on the integration path.

Define the real dilogarithm on the whole real axis by
\begin{equation}
 \mathcal L_2(x)=\operatorname{Re}\operatorname{Li}_2(x+\ii0),
 \qquad
 \frac{\dd\mathcal L_2(x)}{\dd x}=-\frac{\ln|1-x|}{x}.
\end{equation}
For $x>1$ it can be evaluated without a complex branch choice as
\begin{equation}
 \mathcal L_2(x)=\frac{\pi^2}{6}-\ln x\ln(x-1)
 -\operatorname{Li}_2(1-x),\qquad x>1.
\end{equation}
This identity follows by differentiating both sides and matching their
continuous limits at $x=1$. For $x\leq1$ the ordinary real dilogarithm
applies. Thus an argument larger than one does not introduce an imaginary
contribution to $K$. For $g$ and $h$ drawn from the $h_j$, put
\begin{align}
 \mathcal J(Y;g,h)
 &=\ln|1-g/h|\ln|1-Y/g|
 -\mathcal L_2\!\left(\frac{Y-g}{h-g}\right)
 +\mathcal L_2\!\left(-\frac{g}{h-g}\right),
 \qquad g\ne h,\\
 \mathcal J(Y;g,g)&=\frac12\ln^2|1-Y/g|.
\end{align}
Differentiation establishes the required primitive directly,
\begin{equation}
 \frac{\partial\mathcal J(Y;g,h)}{\partial Y}
 =\frac{\ln(1-Y/h)}{Y-g},\qquad \mathcal J(0;g,h)=0.
\end{equation}
Indeed, for $g\ne h$ the two logarithms produced by differentiation
combine to $\ln|(h-Y)/h|$. Positivity proved above removes the absolute
value from this last expression. The closed result is therefore
\begin{align}
 W_j&=L_0\ln|1-Y/h_j|
 +\sum_{k=1}^4\epsilon_k\mathcal J(Y;h_j,h_k),\\
 K&=\boxed{\frac{z_e}{aE}\sum_{j=1}^4 A_jW_j
 -\frac{\operatorname{artanh}\beta}{\beta}.}
 \label{eq:comp-soft-K-dilog}
\end{align}
This is a finite sum of logarithms and dilogarithms with explicit
arguments. No parameter quadrature remains in the finite soft function.

\paragraph{Angular endpoints.}
At $t=0$, $M_x=1$ and $v_x=\beta$ throughout the interval, which gives
\begin{equation}
 K(s,0)=0.
\end{equation}
At $t=-4q^2$ the rationalizing variable used above degenerates because
$\kappa=0$. The endpoint can instead be integrated directly. Set
$w=\operatorname{artanh}\beta$. Then
\begin{align}
 I\big|_{t=-4q^2}
 &=\frac1{\beta E^2}\int_0^\beta
 \frac{\operatorname{artanh}v}{v(1-v^2)}\,\dd v
 =\frac{J_\beta}{\beta E^2},\\
 J_\beta&=\frac{\operatorname{Li}_2(\beta)-\operatorname{Li}_2(-\beta)}2
 +w\ln(2\cosh w)-\frac{w^2}2
 -\frac12\operatorname{Li}_2(-e^{-2w})-\frac{\pi^2}{24},\\
 K(s,-4q^2)&=\frac{1+2q^2}{\beta E^2}J_\beta-\frac w\beta.
 \label{eq:comp-soft-backward-endpoint}
\end{align}
To obtain the second line, split $1/[v(1-v^2)]$ into
$1/v+v/(1-v^2)$. The first term integrates to the difference of two
dilogarithms. In the second term $v=\tanh u$ gives
$\int_0^w u\tanh u\,\dd u$, which is integrated by parts using
$\ln\cosh u=u+\ln(1+e^{-2u})-\ln2$.
All dilogarithm arguments in this endpoint expression are less than one.

\subsection{The integrated infrared coefficient}
The coefficient multiplying the soft logarithm can also be integrated
analytically over the elastic angle. Define the spin-averaged Born
squared amplitude with its factor $e_{\mathrm T}^4$ removed by $\mathcal B$.
In terms of $x=s-1$ and $y=1-u=x+t$, its invariant expression is
\begin{equation}
 \mathcal B=2\left[\frac xy+\frac yx
 -4\left(\frac1y-\frac1x\right)
 +4\left(\frac1y-\frac1x\right)^2\right].
\end{equation}
The forward limit is $\mathcal B=4$. At threshold this expression tends
to $2(1+c^2)$, which also checks the normalization. Write the required
angular coefficient as
\begin{equation}
 \mathcal C_{\mathrm{IR}}(s)=\int_{-1}^1
 \mathcal B(s,t)\,[b_e(t)-1]\,\dd c,
 \qquad b_e(t)=\frac{z_e\operatorname{arcosh}z_e}{\sqrt{z_e^2-1}}.
\end{equation}
The rapidity variable $r=\exp(\operatorname{arcosh}z_e)$ removes the
square root in $b_e$. Reversing the integration limits gives
\begin{align}
 t&=2-r-r^{-1},\qquad y=s+1-r-r^{-1},\qquad 1\leq r\leq s,\\
 \mathcal C_{\mathrm{IR}}(s)
 &=\frac{2s}{x^2}\int_1^s
 \left[R_+(r)\ln r-R_-(r)\right]\dd r,\\
 R_\eta(r)&=\mathcal B(s,2-r-r^{-1})(1+\eta r^{-2}),
 \qquad \eta=\pm1.
\end{align}
The last integral is a rational function times a single logarithm.
For completeness all its partial-fraction coefficients can be stated
explicitly. Put
\begin{align}
 A&=x-4-\frac8x,\qquad B_0=\frac2x\left(s+5+\frac4x\right),\\
 z_\pm&=\frac{s+1\pm\sqrt{(s-1)(s+3)}}2,\\
 d_z&=-1+z^{-2},\qquad e_z=-2z^{-3},
 \qquad f_{\eta z}=1+\eta z^{-2}.
\end{align}
Then
\begin{align}
 R_\eta(r)&=\sum_{n=-3}^{1}p_{\eta n}r^n
 +\sum_{z=z_\pm}\left[\frac{c_{\eta z}}{r-z}
 +\frac{d_{\eta z}}{(r-z)^2}\right],\\
 p_{\eta,1}&=-\frac2x,\quad p_{\eta,0}=B_0,
 \quad p_{\eta,-3}=-\frac{2\eta}x,\quad
 p_{\eta,-2}=\eta B_0,\\
 p_{\eta,-1}&=-\frac2x+\eta\left(-\frac2x-2A\right),\\
 d_{\eta z}&=\frac{8f_{\eta z}}{d_z^2},\qquad
 c_{\eta z}=f_{\eta z}\left(\frac{2A}{d_z}
 -\frac{8e_z}{d_z^3}\right)-\frac{16\eta}{z^3d_z^2}.
\end{align}
The coefficients follow by expanding $y$ to second order at each of its
zeros. Both zeros are outside the interval: $0<z_-<1<s<z_+$.
No pole displacement or continuation is required for this angular integral.

Define elementary power primitives and the two pole primitives by
\begin{align}
 U_n(r)&=\frac{r^{n+1}}{n+1},\qquad
 V_n(r)=r^{n+1}\left[\frac{\ln r}{n+1}-\frac1{(n+1)^2}\right],
 \qquad n\ne-1,\\
 U_{-1}(r)&=\ln r,\qquad V_{-1}(r)=\frac12\ln^2r,\\
 H_1(r,z)&=\ln|r-z|,\qquad H_2(r,z)=-\frac1{r-z},\\
 J_1(r,z)&=\ln r\ln|1-r/z|+\mathcal L_2(r/z),\\
 J_2(r,z)&=-\frac{\ln r}{r-z}
 +\frac{\ln|r-z|-\ln r}{z}.
\end{align}
Their derivatives are, respectively, $r^n$, $r^n\ln r$,
$(r-z)^{-j}$, and $(r-z)^{-j}\ln r$. The integrated result is
\begin{align}
 \mathcal F_+(r)&=\sum_{n=-3}^1p_{+,n}V_n(r)
 +\sum_{z=z_\pm}[c_{+z}J_1(r,z)+d_{+z}J_2(r,z)],\\
 \mathcal F_-(r)&=\sum_{n=-3}^1p_{-,n}U_n(r)
 +\sum_{z=z_\pm}[c_{-z}H_1(r,z)+d_{-z}H_2(r,z)],\\
 \Delta\mathcal F_\eta&=\mathcal F_\eta(s)-\mathcal F_\eta(1),\\
 \mathcal C_{\mathrm{IR}}(s)&=\boxed{\frac{2s}{(s-1)^2}
 [\Delta\mathcal F_+-\Delta\mathcal F_-].}
 \label{eq:comp-soft-integrated-ir}
\end{align}
This expression contains no angular integral. The cancellation of the total soft logarithm can therefore be checked
with its coefficient already integrated analytically.

\subsection{Exchanging the finite soft angular integrals}
\label{sec:comp-soft-angular-exchange}
The remaining finite soft angular coefficient is
\begin{equation}
 \mathcal C_K(s)=\int_{-1}^1\mathcal B(s,t)K(s,t)\,\dd c.
\end{equation}
Rather than integrate the individual dilogarithms in
\eqref{eq:comp-soft-K-dilog}, it is useful to exchange the original
parameter and angular integrations. Put $T=-t$, $h=1+T/4$, and
\begin{equation}
 f(M)=\frac{\operatorname{artanh}\sqrt{1-M/E^2}}
 {M\sqrt{1-M/E^2}},\qquad
 I(T)=\frac1{2\sqrt{h-1}}\int_1^h\frac{f(M)}{\sqrt{h-M}}\,\dd M.
\end{equation}
This is obtained from $M=h-(h-1)y^2$ in the original parameter integral.
Since $1\leq h\leq E^2$, interchanging the two integrations yields
\begin{equation}
 \int_0^{4q^2}\mathcal B(T)(1+T/2)I(T)\,\dd T
 =2\int_1^{E^2}f(M)\,\dd M
 \int_M^{E^2}\frac{\mathcal B(h)(2h-1)}
 {\sqrt{(h-1)(h-M)}}\,\dd h.
\end{equation}
All endpoint singularities here are integrable. The substitution
$u=\sqrt{(h-M)/(h-1)}$ makes the inner integral rational. Its upper
limit is $U=\sqrt{E^2-M}/q$. Define
\begin{equation}
 k^2=\frac{1+xU^2}{s},\qquad x=s-1,
 \qquad h=\frac{M-u^2}{1-u^2}.
\end{equation}
The rational function that occurs is
\begin{equation}
 \frac{\mathcal B(h)(2h-1)}{1-u^2}
 =\frac{a_1}{1-u^2}+\frac{a_2}{(1-u^2)^2}
 +\frac{a_3}{(1-u^2)^3}
 +\frac{b_1}{k^2-u^2}+\frac{b_2}{(k^2-u^2)^2},
\end{equation}
with explicit coefficients
\begin{align}
 a_1&=-\frac{x^3-6x^2-16x-8}{x^2},
 &a_2&=\frac{(1-k^2)(x^2+2x+4)}x,\\
 a_3&=-x(1-k^2)^2,
 &b_1&=\frac{x^3-2x^2-16x-8}{x^2},\\
 b_2&=\frac{4(1-k^2)(x+2)}{x^2}.
\end{align}
The elementary primitives needed at the upper limit are
\begin{align}
 F_1(U;k)&=\frac{\operatorname{artanh}(U/k)}k,\\
 F_2(U;k)&=\frac{U}{2k^2(k^2-U^2)}+\frac{F_1(U;k)}{2k^2},\\
 F_3(U;k)&=\frac{U}{4k^2(k^2-U^2)^2}+\frac{3F_2(U;k)}{4k^2}.
\end{align}
They vanish at $U=0$ and differentiate to $(k^2-U^2)^{-n}$ when $k$
is held fixed. In the inner integration it is indeed fixed. Consequently
\begin{align}
 H(U)&=a_1F_1(U;1)+a_2F_2(U;1)+a_3F_3(U;1)
 +b_1F_1(U;k)+b_2F_2(U;k),\\
 \mathcal C_K(s)&=\frac4\beta\int_0^1
 \frac{\operatorname{artanh}(\beta U)}{E^2(1-\beta^2U^2)}
 H(U)\,\dd U-\frac{\operatorname{artanh}\beta}{\beta}
 \mathcal C_0(s),
 \label{eq:comp-soft-K-angular-reduced}\\
 \mathcal C_0(s)&=\int_{-1}^1\mathcal B\,\dd c
 =\frac{4s}{x^2}[P(x)-P(x/s)],\\
 P(y)&=\frac{y^2}{2x}+\left(\frac4x+\frac4{x^2}\right)y
 +\left(x-4-\frac8x\right)\ln y-\frac4y.
\end{align}
The square root $k$ is positive, and
$k^2-U^2=(1-U^2)/s>0$ for $0\leq U<1$. Both inverse hyperbolic tangents
in $H$ are therefore real. This reduction integrates the original elastic
angle analytically before performing the remaining soft parameter integral.

The one remaining square root can itself be rationalized by
\begin{align}
 w&=\sqrt{x}U+\sqrt{1+xU^2},
 &1\leq w&\leq\sqrt{x}+\sqrt{s},\\
 U&=\frac{w-w^{-1}}{2\sqrt{x}},
 &k&=\frac{w+w^{-1}}{2\sqrt{s}}.
\end{align}
Thus \eqref{eq:comp-soft-K-angular-reduced} becomes an integral of
rational functions of $w$ multiplying one or two logarithms of rational
functions. It requires at most weight-three multiple polylogarithms.
The endpoints and all logarithm arguments have been obtained directly
from the physical interval, without selecting branches by comparison
with a numerical answer.

\paragraph{Reduction to a finite polylogarithmic expression.}
It is useful to state an algebraic reduction which avoids printing a
large list of weight-three terms. It terminates after at most two
reductions for the expression below and contains no quadrature.
Set $z=w-1$, $Z=\sqrt{s}+\sqrt{x}-1$, and regard $U$ and $k$ as the
rational functions of $w=1+z$ just obtained. Write
\begin{align}
 H&=h_0+h_1\operatorname{artanh}U
       +h_k\operatorname{artanh}(U/k),\\
 h_0&=\frac{a_2U}{2(1-U^2)}
 +a_3\left[\frac{U}{4(1-U^2)^2}
 +\frac{3U}{8(1-U^2)}\right]
 +\frac{b_2U}{2k^2(k^2-U^2)},\\
 h_1&=a_1+\frac{a_2}2+\frac{3a_3}8,
 \qquad h_k=\frac{b_1}k+\frac{b_2}{2k^3},\\
 R_j(z)&=\frac{4U'(z)h_j}{\beta E^2(1-\beta^2U^2)},
 \qquad j=0,1,k.
\end{align}
Define the positive root
\begin{equation}
 r=\sqrt{1+(\sqrt{x}/\beta)^2}-\sqrt{x}/\beta,\qquad 0<r<1.
\end{equation}
Each $R_j$ is explicitly rational. Its poles in the $w$ plane belong to
\begin{equation}
 w=0,\quad \pm\ii,\quad \pm r,\quad \pm r^{-1},
\end{equation}
The orders at $w=0$ are three for $R_0$ and
four for $R_1$, and $R_k$ has no pole there. At $w=\pm\ii$, $R_0$ has
simple poles and $R_k$ has double poles. The four real poles are simple.
The polynomial parts of $R_0$, $R_1$, and $R_k$ have degrees one, two,
and no polynomial part, respectively. These statements follow by
substituting the displayed rational expressions for $U$ and $k$.
Apparent endpoint poles in $h_0$
cancel because $1-k^2=x(1-U^2)/s$ and $k^2-U^2=(1-U^2)/s$.

\paragraph{Reducing the kernel before integration.}
Put $X=x=s-1$ in this paragraph. An integration by parts removes the
poles at $w=\pm\ii$ before any polylogarithmic primitive is constructed.
It combines the elementary and root-dependent parts of the original
soft kernel, so no integrated cross section is needed.

Keep $X=s-1$, $\beta=X/(X+2)$ and $E^2=(X+2)^2/(4s)$, and define
\begin{equation}
 k^2=\frac{1+XU^2}{s},\qquad d=1-\beta^2U^2,\qquad
 a=\operatorname{artanh}(\beta U),\qquad
 f=\operatorname{artanh}(U/k).
\end{equation}
In this section $d$ denotes this rational function, rather than the
spacetime dimension.  On $0<U<1$, all square roots are positive and
both functions are real.  Separate the elementary part of $H$ as
\begin{align}
 h_0&=\overline h_0+\frac{2(X+2)U}{Xk^2},\\
 \overline h_0&=U\left[\frac{X^2+2X+4}{2s}
 -\frac{X^3}{4s^2}-\frac{3X^3(1-U^2)}{8s^2}\right],\\
 h_k&=\frac{b_1}{k}+\frac{2(X+2)(1-U^2)}{Xsk^3},\qquad
 b_1=\frac{X^3-2X^2-16X-8}{X^2}.
\end{align}
The identities
\begin{equation}
 \frac{\dd f}{\dd U}=\frac1{k(1-U^2)},\qquad
 p(U)=\frac{2(X+2)U(1-U^2)}X
\end{equation}
give the exact integration-by-parts relation
\begin{equation}
 \begin{aligned}
 \frac a d\left[\frac{2(X+2)U}{Xk^2}+h_kf\right]
 ={}&\frac{\dd}{\dd U}\left(\frac{apf}{dk}\right)\\
 &+\frac f k\left[
 a\left(\frac{b_1}d+\frac{16sU^2}{X(X+2)d^2}\right)
 -\frac{2U(1-U^2)}{d^2}\right].
 \end{aligned}
 \label{eq:comp-ck-pole-removal}
\end{equation}
Its boundary term vanishes at both ends.  At $U=0$, $a$ and $f$
vanish linearly.  At $U\to1$, $p=O(1-U)$ while $f$ grows only as
$\ln(1-U)$, and $d\to1-\beta^2>0$.  There is therefore no boundary
constant to supply from a different calculation.

Use the previous rational variable
$U=(w-w^{-1})/(2\sqrt X)$.  The remaining root factor cancels from
its differential,
\begin{equation}
 \frac{\dd U}{k}=\sqrt{\frac sX}\frac{\dd w}{w}.
\end{equation}
Writing $z=w-1$, define four rational kernels
\begin{align}
 \widetilde R_0&=\frac4{\beta E^2}\frac{U'\overline h_0}{d},&
 \widetilde R_1&=\frac4{\beta E^2}\frac{U'h_1}{d},\\
 \widetilde R_k&=\frac4{\beta E^2}\sqrt{\frac sX}\frac1w
 \left[\frac{b_1}d+\frac{16sU^2}{X(X+2)d^2}\right],&
 \widetilde R_f&=-\frac8{\beta E^2}\sqrt{\frac sX}
 \frac{U(1-U^2)}{wd^2}.
\end{align}
Here $U'=\dd U/\dd z$ and
\begin{equation}
 h_1=-\frac{X^3-6X^2-16X-8}{X^2}
 +\frac{(1-k^2)(X^2+2X+4)}{2X}
 -\frac{3X}{8}(1-k^2)^2.
\end{equation}
The logarithms are $L_\beta=\operatorname{artanh}(\beta U)$,
$L_1=\operatorname{artanh}U$, and $L_k=\operatorname{artanh}(U/k)$.
The transformed integrand is
\begin{equation}
 \widetilde R_0 L_\beta+\widetilde R_1L_\beta L_1
 +\widetilde R_kL_\beta L_k+\widetilde R_f L_k.
\end{equation}
Every denominator now factors into powers of $w$ and
\begin{equation}
 Q_w=-Xw^4+(4X^2+18X+16)w^2-X.
\end{equation}
The roots of $Q_w$ are the four real poles
$\pm r,\pm r^{-1}$, where
$r=\sqrt{1+X/\beta^2}-\sqrt X/\beta$.
The factor $1+w^2$ is absent.
Thus the entire GPL expression, including the letters introduced by
partial fractions, has only real letters.


For a finite list of pairs $(a,\ell_a)$ write
$L(z)=\sum_a\ell_aG(a;z)$, where $G$ is the usual multiple
polylogarithm with kernels $1/(z-a)$. The three lists needed here are
\begin{align}
 \mathcal A_\beta&=\tfrac12\{(r-1,+1),(-r^{-1}-1,+1),
 (r^{-1}-1,-1),(-r-1,-1)\},\\
 \mathcal A_1&=\tfrac12\{(W^{-1}-1,+1),(-W-1,+1),
 (W-1,-1),(-W^{-1}-1,-1)\},\\
 \mathcal A_k&=\tfrac12\{(W^{-1}-1,+1),(-W^{-1}-1,+1),
 (W-1,-1),(-W-1,-1)\},\\
 W&=\sqrt{s}+\sqrt{x}.
\end{align}
Here the factor $1/2$ multiplies the weights, not the letters.
Factoring $1\pm\beta U$, $1\pm U$, and $k\pm U$ proves
\begin{equation}
 L_\beta(z)=\operatorname{artanh}(\beta U),\qquad
 L_1(z)=\operatorname{artanh}U,\qquad
 L_k(z)=\operatorname{artanh}(U/k).
\end{equation}
All three vanish at $z=0$, which fixes the constants in the factorized
logarithms. Their real letters lie outside $0<z<Z$.
The two functions with endpoint logarithms are defined by their limits
from $z<Z$.

Here is an explicit algebraic rule for removing the rational factors.
Let $\mathfrak T[R;\boldsymbol a](z)$ be linear in the rational function
$R$. Expand $R$ into its polynomial part and partial fractions. For a pole
$a$ of order $m_a$, the coefficient of $(z-a)^{-n}$ is explicitly
\begin{equation}
 C_{a,n}=\left.\frac1{(m_a-n)!}
 \frac{\dd^{m_a-n}}{\dd z^{m_a-n}}
 [(z-a)^{m_a}R(z)]\right|_{z=a},\qquad 1\leq n\leq m_a.
\end{equation}
The polynomial part is obtained by polynomial division. The known pole
list above makes all these operations finite. For a
nonempty word $\boldsymbol a=(b,\boldsymbol c)$ use
\begin{align}
 \mathfrak T[(z-a)^{-1};\boldsymbol a](z)
 &=G(a,\boldsymbol a;z),\\
 \mathfrak T[(z-a)^{-n};b,\boldsymbol c](z)
 &=-\frac{G(b,\boldsymbol c;z)}{(n-1)(z-a)^{n-1}}
 +\frac1{n-1}\mathfrak T\!\left[
 \frac1{(z-a)^{n-1}(z-b)};\boldsymbol c\right](z),
 \quad n\geq2,\\
 \mathfrak T[z^n;b,\boldsymbol c](z)
 &=\frac{z^{n+1}}{n+1}G(b,\boldsymbol c;z)
 -\frac1{n+1}\mathfrak T\!\left[
 \frac{z^{n+1}}{z-b};\boldsymbol c\right](z),
 \quad n\geq0.
\end{align}
For the empty word, use the elementary rules
\begin{align}
 \mathfrak T[z^n;\varnothing](z)&=\frac{z^{n+1}}{n+1},\\
 \mathfrak T[(z-a)^{-1};\varnothing](z)&=G(a;z),\\
 \mathfrak T[(z-a)^{-n};\varnothing](z)
 &=\frac{(z-a)^{1-n}-(-a)^{1-n}}{1-n},\qquad n\geq2.
\end{align}
Every pole appearing here has $a\ne0$. Differentiation proves that
$\mathfrak T[R;\boldsymbol a]'=R\,G(\boldsymbol a;z)$ and its value
at zero is zero. For repeated poles and polynomial terms the word
length decreases. Simple poles terminate immediately. Hence these
rules specify a finite combination of rational functions and ordinary
multiple polylogarithms, rather than a new unevaluated integral.

The endpoint expression for the finite soft angular coefficient is
\begin{align}
 \mathcal S_0(z)&=\sum_{a\in\mathcal A_\beta}
 \ell_a\mathfrak T[R_0;a](z),\\
 \mathcal S_j(z)&=\sum_{a\in\mathcal A_\beta}
 \sum_{b\in\mathcal A_j}\ell_a\ell_b
 \left\{\mathfrak T[R_j;a,b](z)
 +\mathfrak T[R_j;b,a](z)\right\},\qquad j=1,k,\\
 \mathcal C_K(s)&=\lim_{z\to Z^-}
 [\mathcal S_0(z)+\mathcal S_1(z)+\mathcal S_k(z)]
 -\frac{\operatorname{artanh}\beta}{\beta}\mathcal C_0(s).
 \label{eq:comp-soft-K-angular-GPL}
\end{align}
The two words in $\mathcal S_j$ are the shuffle identity for the product
of two weight-one logarithms. Formula \eqref{eq:comp-soft-K-angular-GPL}
is an algebraically specified weight-three polylogarithmic expression.
An expanded list of its coefficients can be generated using only
partial fractions and the finite rules above. The limit is finite:
the original $U$ integrand has at most an integrable logarithm at $U=1$.

\paragraph{The reduced soft endpoint.}
For an explicit closed endpoint expression, apply the finite
primitive operator $\mathfrak T$ just defined to the four reduced
kernels and the logarithm lists $\mathcal A_\beta,\mathcal A_1,\mathcal A_k$. The notation $\mathfrak T[R;L_aL_b]$ means the finite
sum over their weights after applying
$G(a)G(b)=G(a,b)+G(b,a)$.  Put
\begin{equation}
 \begin{aligned}
 \widetilde\Psi(z)={}&
 \mathfrak T[\widetilde R_0;L_\beta](z)
 +\mathfrak T[\widetilde R_1;L_\beta L_1](z)\\
 &+\mathfrak T[\widetilde R_k;L_\beta L_k](z)
 +\mathfrak T[\widetilde R_f;L_k](z).
 \end{aligned}
\end{equation}
Each primitive is normalized to vanish at $z=0$.  The simplified
answer is
\begin{equation}
 \boxed{\mathcal C_K(s)=
 \lim_{z\to Z^-}\widetilde\Psi(z)
 -\frac{\operatorname{artanh}\beta}{\beta}\mathcal C_0(s),
 \qquad Z=\sqrt s+\sqrt X-1.}
 \label{eq:comp-ck-simplified-endpoint}
\end{equation}
The combined endpoint is finite.  All terms are GPLs of weight at
most three with algebraic endpoint coefficients.  There are two further exact residue cancellations.  If
$P_1$ denotes the primitive of the higher-order pole part of
$\widetilde R_1$ at $w=0$, then
\begin{equation}
 \operatorname*{Res}_{w=0}\widetilde R_1=0,\qquad
 \operatorname*{Res}_{w=0}\left[
 \widetilde R_0-P_1\frac{\dd L_1}{\dd w}\right]=0.
\end{equation}
The first identity removes 32 weight-three words introduced by a
vanishing simple-pole coefficient.  The second removes four
weight-two words.  They follow by substituting the displayed rational
kernels, so no small numerical coefficient is discarded.
After these cancellations, the finite rule produces 190 terms instead
of 306, and removes both imaginary letters.  The same physical interval fixes all logarithms.


\paragraph{The finite endpoint limit.} The convergence of the endpoint
expression follows from its local expansion.
Let $d=1-U$ and let $\mathcal I(U)$ denote the integrand including the
factor $4/\beta$ in \eqref{eq:comp-soft-K-angular-reduced}. Define
\begin{align}
 h_{0*}&=\frac{x^2+2x+4}{2s}-\frac{x^3}{4s^2}+\frac{2(x+2)}x,\\
 B_*&=h_{0*}+2\ln2+\frac{b_1}2\ln s.
\end{align}
The rational coefficients have finite limits at $U=1$, and
$h_1+h_k=4$ there. Also $1-U/k=d/s+O(d^2)$. Hence
\begin{align}
 \mathcal I(1-d)&=\frac{4\operatorname{artanh}\beta}{\beta}
 [-2\ln d+B_*]+O(d\ln d),\\
 \int_{1-d}^1\mathcal I(U)\,\dd U
 &=\frac{4\operatorname{artanh}\beta}{\beta}
 d[-2\ln d+2+B_*]+O(d^2\ln d).
\end{align}
Since $d\ln d\to0$, the omitted tail vanishes. This establishes the
existence of the combined polylogarithmic endpoint limit.

\subsection{Changing the soft energy frame}
A useful independent check is to place the soft-photon energy bound in
the final electron rest frame. Write $z_e=\cosh\eta$ and
$v=\tanh\eta$. In this frame the final electron has energy one, the
initial electron has energy $z_e$, and their interpolated momentum has
\begin{equation}
 E_x=1+x(z_e-1),\qquad
 M_x=1+2(z_e-1)x(1-x),\qquad
 v_x=\frac{x\sqrt{z_e^2-1}}{1+x(z_e-1)}.
\end{equation}
In this paragraph $x$ is a Feynman parameter. The finite cross term
integrates directly after changing from $x$ to $v_x$,
\begin{equation}
 \int_0^1\frac{\operatorname{artanh}v_x}{v_xM_x}\,\dd x
 =\frac{J(v)}{\sqrt{z_e^2-1}},\qquad
 J(v)=\int_0^v\frac{\operatorname{artanh}u}{u(1-u^2)}\,\dd u.
\end{equation}
The integral $J$ is the dilogarithmic expression $J_\beta$ above with
$\beta$ replaced by $v$ and $w$ by $\eta$. Including the two self terms
gives the closed finite function
\begin{equation}
 K_{\mathrm{rest}}(t)=\frac{J(v)}v-\frac{b_e+1}2,
 \qquad b_e=\frac\eta v.
\end{equation}
Thus the difference between a rest-frame energy bound $\sqrt{s}\varepsilon$
and a center-of-mass energy bound $\varepsilon$ is
\begin{equation}
 \frac{\delta_{\mathrm{rest}}-\delta_{\mathrm{CM}}}{\alpha_{\mathrm T}}
 =\frac2\pi\left[(b_e-1)\ln\sqrt{s}+K-K_{\mathrm{rest}}\right].
 \label{eq:comp-soft-frame-change}
\end{equation}
This finite change also follows directly from the angular integral of
the eikonal factor multiplied by $\ln[\sqrt{s}/(p'\cdot n)]$, where
$n=(1,\widehat{\boldsymbol k})$ in the center-of-mass frame.

For example, consider the asymmetric invariant cuts
$2P\cdot k_1>\delta$ and $2p'\cdot k_2>\delta$, with
$P=p+k$ and $\delta=2\sqrt{s}\varepsilon$.
Since $2P\cdot k_1+2p'\cdot k_2=s-1$, their omitted soft regions do not
overlap for sufficiently small $\delta$. The two soft terms have weights
$1/2$ from the identical-photon phase-space factor. Their average differs
from the original center-of-mass soft term by half of
\eqref{eq:comp-soft-frame-change}. In the total coefficient multiplying
$\alpha_{\mathrm T}^3/m^2$, the shift is
\begin{equation}
 \Delta C_{\mathrm{VS}}=
 \frac1{2s}\left[\mathcal C_K-\mathcal C_{K,\mathrm{rest}}
 +\frac{\ln s}{2}\mathcal C_{\mathrm{IR}}\right].
\end{equation}
It cannot be omitted if this asymmetric cut is used.

The rest-frame angular coefficient has a short weight-three expression.
Set $r=e^\eta$ and $z=r-1$. Since
\begin{equation}
 \frac{\dd J}{\dd r}=\ln r\left[-\frac1r+\frac1{r-1}+\frac1{r+1}\right],
 \qquad J(0\text{ in }\eta)=0,
\end{equation}
its exact polylogarithmic form is
\begin{equation}
 J=-G(-1,-1;z)+G(0,-1;z)+G(-2,-1;z).
\end{equation}
Let $\widetilde R_+(z)=R_+(1+z)$, with $R_+$ the rational kernel used
in the integrated infrared coefficient. The algebraic map $\mathfrak T$
then gives
\begin{align}
 \mathcal Q(z)&=-\mathfrak T[\widetilde R_+;-1,-1](z)
 +\mathfrak T[\widetilde R_+;0,-1](z)
 +\mathfrak T[\widetilde R_+;-2,-1](z),\\
 \mathcal C_{K,\mathrm{rest}}(s)
 &=\frac{2s}{(s-1)^2}\mathcal Q(s-1)
 -\frac12\mathcal C_{\mathrm{IR}}(s)-\mathcal C_0(s).
\end{align}
This alternative frame is not needed for the symmetric analytic
regulator described next.
